% Options for packages loaded elsewhere
\PassOptionsToPackage{unicode}{hyperref}
\PassOptionsToPackage{hyphens}{url}
\documentclass[
  12pt,
]{ctexart}
\usepackage{xcolor}
\usepackage{amsmath,amssymb}
\setcounter{secnumdepth}{-\maxdimen} % remove section numbering
\usepackage{iftex}
\ifPDFTeX
  \usepackage[T1]{fontenc}
  \usepackage[utf8]{inputenc}
  \usepackage{textcomp} % provide euro and other symbols
\else % if luatex or xetex
  \usepackage{unicode-math} % this also loads fontspec
  \defaultfontfeatures{Scale=MatchLowercase}
  \defaultfontfeatures[\rmfamily]{Ligatures=TeX,Scale=1}
\fi
\usepackage{lmodern}
\ifPDFTeX\else
  % xetex/luatex font selection
\fi
% Use upquote if available, for straight quotes in verbatim environments
\IfFileExists{upquote.sty}{\usepackage{upquote}}{}
\IfFileExists{microtype.sty}{% use microtype if available
  \usepackage[]{microtype}
  \UseMicrotypeSet[protrusion]{basicmath} % disable protrusion for tt fonts
}{}
\makeatletter
\@ifundefined{KOMAClassName}{% if non-KOMA class
  \IfFileExists{parskip.sty}{%
    \usepackage{parskip}
  }{% else
    \setlength{\parindent}{0pt}
    \setlength{\parskip}{6pt plus 2pt minus 1pt}}
}{% if KOMA class
  \KOMAoptions{parskip=half}}
\makeatother
\usepackage{longtable,booktabs,array}
\usepackage{calc} % for calculating minipage widths
% Correct order of tables after \paragraph or \subparagraph
\usepackage{etoolbox}
\makeatletter
\patchcmd\longtable{\par}{\if@noskipsec\mbox{}\fi\par}{}{}
\makeatother
% Allow footnotes in longtable head/foot
\IfFileExists{footnotehyper.sty}{\usepackage{footnotehyper}}{\usepackage{footnote}}
\makesavenoteenv{longtable}
\setlength{\emergencystretch}{3em} % prevent overfull lines
\providecommand{\tightlist}{%
  \setlength{\itemsep}{0pt}\setlength{\parskip}{0pt}}
\usepackage{bookmark}
\IfFileExists{xurl.sty}{\usepackage{xurl}}{} % add URL line breaks if available
\urlstyle{same}
\hypersetup{
  hidelinks,
  pdfcreator={LaTeX via pandoc}}

\author{SCX}
\date{2026}

\begin{document}

\section{SCX 自我进化理论 ---
数学形式化}\label{scx-ux81eaux6211ux8fdbux5316ux7406ux8bba-ux6570ux5b66ux5f62ux5f0fux5316}

\begin{quote}
将 SCX 的自我进化闭环（judge → store → update → re-judge → re-update →
\ldots）形式化为严格的数学理论，并与已知理论框架建立连接。
\end{quote}

\textbf{最后更新}: 2026-06-28

\begin{center}\rule{0.5\linewidth}{0.5pt}\end{center}

\subsection{问题}\label{ux95eeux9898}

SCX 系统在运行中形成以下闭环：

\begin{verbatim}
Judge (S_t 评分) → Store (M_t 记忆库) → Update (S_{t+1} 更新评分) → Re-judge → ...
                                ↘ Update (θ_{t+1} NEP 学生更新) ↗
\end{verbatim}

\textbf{核心问题}：这个闭环是否收敛？收敛到何处？收敛速率如何？在何种条件下发散？

\textbf{困难}：Gatekeeper \(S_t\) 和 NEP 学生 \(f_{\\theta_t}\)
构成耦合动力系统------\(S_t\) 的选择决定了 \(f_{\\theta_t}\)
的训练分布，而 \(f_{\\theta_t}\) 的输出提供反馈给 \(S_t\)。

\begin{center}\rule{0.5\linewidth}{0.5pt}\end{center}

\subsection{核心洞察}\label{ux6838ux5fc3ux6d1eux5bdf}

\begin{verbatim}
SCX 自我进化的正确数学对象是耦合动力系统 (S_t, θ_t, M_t)，
其全局行为由 Lyapunov 函数 V(S, θ) 的单调下降性质主导。
\end{verbatim}

其中 \(S_t: \mathcal{X} \to [0,1]\) 是 gatekeeper 评分函数，\(\theta_t\)
是 NEP 学生参数，\(M_t\) 是单调增长的经验记忆库。

\begin{center}\rule{0.5\linewidth}{0.5pt}\end{center}

\subsection{九个核心文件}\label{ux4e5dux4e2aux6838ux5fc3ux6587ux4ef6}

\begin{longtable}[]{@{}
  >{\raggedright\arraybackslash}p{(\linewidth - 6\tabcolsep) * \real{0.2143}}
  >{\raggedright\arraybackslash}p{(\linewidth - 6\tabcolsep) * \real{0.2143}}
  >{\raggedright\arraybackslash}p{(\linewidth - 6\tabcolsep) * \real{0.3571}}
  >{\raggedright\arraybackslash}p{(\linewidth - 6\tabcolsep) * \real{0.2143}}@{}}
\toprule\noalign{}
\begin{minipage}[b]{\linewidth}\raggedright
文件
\end{minipage} & \begin{minipage}[b]{\linewidth}\raggedright
内容
\end{minipage} & \begin{minipage}[b]{\linewidth}\raggedright
核心结果
\end{minipage} & \begin{minipage}[b]{\linewidth}\raggedright
状态
\end{minipage} \\
\midrule\noalign{}
\endhead
\bottomrule\noalign{}
\endlastfoot
\texttt{01\_symbol\_system.md} & 符号系统与问题设定 & 结构空间
\(\mathbb{X}\)、状态空间 \(\mathbb{Z}\)、更新算子 \(\Phi\)、新假设 SE-A1
至 SE-A6 & \textbf{完整} \\
\texttt{02\_dynamical\_system.md} & 离散动力系统形式化 & Lyapunov 函数
\(V(z_t)\)、不动点存在性（定理 4-5）、单调下降（定理 6）、相图四种区域 &
\textbf{完整} \\
\texttt{03\_online\_learning\_regret.md} & 在线学习 Regret 分析 & Regret
\(\leq 2GR\sqrt{T}\) (OGD)、延迟反馈 Regret
\(\leq 2GR\sqrt{T} + G^2 D_{\max}\sqrt{T}\) (定理 8) & \textbf{完整} \\
\texttt{04\_bayesian\_update.md} & 贝叶斯更新解释 & \(S_t\)
为后验均值、贝叶斯鞅 \(\mathbb{E}[S_{t+1}|\mathcal{F}_t]=S_t\)、Doob
收敛 \(S_t \to S_\infty\) a.s. & \textbf{完整} \\
\texttt{05\_stochastic\_approximation.md} & 随机逼近分析 & Robbins-Monro
形式 \(\theta_{t+1} = \theta_t - \alpha_t \nabla \ell + \xi_t\)、ODE
方法、双时间尺度分析 & \textbf{完整} \\
\texttt{06\_fixed\_point\_convergence.md} &
\textbf{中心定理：不动点与收敛} & \textbf{Theorem SE-1}：有限结构空间 +
Lipschitz + 充分退火 → \((S_t, \theta_t) \to (S^*, \theta^*)\)
a.s.；四种收敛路径 & \textbf{完整} \\
\texttt{07\_completeness.md} & 完备性分析 & \textbf{Theorem
SE-2}：物理约束下存在有限 \(T^*\) 使得 \(t \geq T^*\) 后系统达到
\(\varepsilon\)-近似不动点；Gödel 类比 & \textbf{完整} \\
\texttt{08\_theory\_connections.md} & 与已知理论的连接 & AlphaZero
self-play、贝叶斯优化、主动学习、Solomonoff 归纳的形式化对比表 &
\textbf{完整} \\
\texttt{09\_verification\_report.md} & 验证报告 &
一致性检查、跨定理交叉引用、符号冲突、证明缺口（10 个）、开放问题（5
个） & \textbf{完整} \\
\end{longtable}

\begin{center}\rule{0.5\linewidth}{0.5pt}\end{center}

\subsection{核心定理}\label{ux6838ux5fc3ux5b9aux7406}

\subsubsection{Theorem SE-1: Convergence of SCX
Self-Evolution}\label{theorem-se-1-convergence-of-scx-self-evolution}

\textbf{条件 (C1-C7)}： 1. 结构空间 \(\mathbb{X}\)
有限（有限数据、有限精度） 2. NEP 学生 \(f_\theta\) 关于 \(\theta\) 是
\(L_f\)-Lipschitz 的 3. Gatekeeper \(S_t\) 关于其参数是
\(L_S\)-Lipschitz 的 4. 学习率满足 Robbins-Monro
条件：\(\sum \alpha_t = \infty, \sum \alpha_t^2 < \infty\) 5. 给定
\(S_t\) 的条件下，记忆库采样是 i.i.d. 的 6.
充分退火：\(\alpha_t \to 0\)，且 gatekeeper 更新频率 \(\to 0\) 7.
Gatekeeper
更新步长有界：\(\|S_{t+1} - S_t\|_\infty \leq \eta_t\)，\(\sum \eta_t < \infty\)

\textbf{结论}：序列 \((S_t, \theta_t)\) 几乎必然收敛到联合不动点
\((S^*, \theta^*)\)，其中： - \(S^*\)
自洽：\(S^*(x) = \mathbb{P}(\text{correct} \mid x, \mathcal{M}_\infty)\)
- \(\theta^*\) 是 \(S^*\) 诱导分布下 NEP 期望损失的局部极小值

\textbf{证明路径}：Lemma SE-1.1 (Lyapunov 超鞅) + Lemma SE-1.2
(有限类型) + Lemma SE-1.3 (步长消失) + Lemma SE-1.4 (极限点即不动点)。

\subsubsection{Theorem SE-2: Completeness
Bound}\label{theorem-se-2-completeness-bound}

在物理约束（有限数据 \(N_{\max}\)、有限精度
\(\varepsilon_{\text{mach}}\)、有限计算）下，存在有限时间 \(T^*\)
使得对所有 \(t \geq T^*\)，系统处于 \(\varepsilon\)-近似不动点。

\begin{center}\rule{0.5\linewidth}{0.5pt}\end{center}

\subsection{四种收敛路径}\label{ux56dbux79cdux6536ux655bux8defux5f84}

\begin{longtable}[]{@{}
  >{\raggedright\arraybackslash}p{(\linewidth - 4\tabcolsep) * \real{0.3333}}
  >{\raggedright\arraybackslash}p{(\linewidth - 4\tabcolsep) * \real{0.3333}}
  >{\raggedright\arraybackslash}p{(\linewidth - 4\tabcolsep) * \real{0.3333}}@{}}
\toprule\noalign{}
\begin{minipage}[b]{\linewidth}\raggedright
路径
\end{minipage} & \begin{minipage}[b]{\linewidth}\raggedright
特征
\end{minipage} & \begin{minipage}[b]{\linewidth}\raggedright
条件
\end{minipage} \\
\midrule\noalign{}
\endhead
\bottomrule\noalign{}
\endlastfoot
\textbf{I. 经典收敛} & \((S_t, \theta_t) \to (S^*, \theta^*)\)，单调改进
& C1-C7 全部满足，充分退火 \\
\textbf{II. 极限环} & 系统在有限配置间周期振荡 & 退火不充分，耦合过强 \\
\textbf{III. 永动发现} & 新结构持续发现，\(\mathcal{M}_t\) 无限增长 &
开放物理世界，无限探索预算 \\
\textbf{IV. 发散崩溃} & \(S_t\) 退化，质量下降 & 反馈回路断裂，NEP
反馈噪声过大 \\
\end{longtable}

路径 II-IV 的精确分界条件是\textbf{开放问题}（见
\texttt{09\_verification\_report.md}）。

\begin{center}\rule{0.5\linewidth}{0.5pt}\end{center}

\subsection{与现有理论体系的关系}\label{ux4e0eux73b0ux6709ux7406ux8bbaux4f53ux7cfbux7684ux5173ux7cfb}

\subsubsection{从属关系}\label{ux4eceux5c5eux5173ux7cfb}

\begin{verbatim}
Theorem 3 (不可识别性) ─── 提供自进化必要性论证 ──→ 自进化闭环必须存在
        │
        ▼
Theorem 1 (噪声检测) ──── 为 S_t 初始化提供保证 ──→ Theorem SE-1 的初始条件
        │
        ▼
Theorem SE-1 (自进化收敛) ─── 中心结果 ──→ 闭环渐近收敛到自洽点
        │
        ├──→ Theorem 2 (弱特征界) ── 限制改进速度
        ├──→ Theorem 4' (最小最大最优) ── 提供 S^* 的渐近最优性
        └──→ Theorem SE-2 (完备性界) ── 物理有限性强制有限时间收敛
\end{verbatim}

\subsubsection{新引入假设 (SE-A1 至
SE-A6)}\label{ux65b0ux5f15ux5165ux5047ux8bbe-se-a1-ux81f3-se-a6}

\begin{longtable}[]{@{}lll@{}}
\toprule\noalign{}
假设 & 内容 & 对应现有假设 \\
\midrule\noalign{}
\endhead
\bottomrule\noalign{}
\endlastfoot
\textbf{SE-A1} & Lyapunov 下降 & \emph{新引入，核心假设} \\
\textbf{SE-A2} & 有限结构空间 & 物理约束（非统计学假设） \\
\textbf{SE-A3} & Lipschitz NEP & 平滑性假设 \\
\textbf{SE-A4} & 条件 i.i.d. 采样 & A1-A2 的扩展 \\
\textbf{SE-A5} & Robbins-Monro 学习率 & 标准 SA 条件 \\
\textbf{SE-A6} & Gatekeeper 更新有界 & 退火条件 \\
\end{longtable}

\begin{center}\rule{0.5\linewidth}{0.5pt}\end{center}

\subsection{与已知理论的连接摘要}\label{ux4e0eux5df2ux77e5ux7406ux8bbaux7684ux8fdeux63a5ux6458ux8981}

\begin{longtable}[]{@{}
  >{\raggedright\arraybackslash}p{(\linewidth - 4\tabcolsep) * \real{0.2273}}
  >{\raggedright\arraybackslash}p{(\linewidth - 4\tabcolsep) * \real{0.2273}}
  >{\raggedright\arraybackslash}p{(\linewidth - 4\tabcolsep) * \real{0.5455}}@{}}
\toprule\noalign{}
\begin{minipage}[b]{\linewidth}\raggedright
理论框架
\end{minipage} & \begin{minipage}[b]{\linewidth}\raggedright
核心机制
\end{minipage} & \begin{minipage}[b]{\linewidth}\raggedright
与 SCX 自进化的关键差异
\end{minipage} \\
\midrule\noalign{}
\endhead
\bottomrule\noalign{}
\endlastfoot
\textbf{AlphaZero} & Policy iteration + MCTS & AlphaZero
生成合成游戏；SCX 依赖真实 NEP 数据 \\
\textbf{贝叶斯优化} & Acquisition + Surrogate & BO 优化静态黑箱；SCX
的目标分布在演化 \\
\textbf{主动学习} & Query strategy + Oracle & AL 的 oracle 即时反馈；SCX
的 NEP 延迟反馈 \\
\textbf{Solomonoff 归纳} & Universal prior & Solomonoff 不可计算；SCX
限制在有限假设类 \\
\end{longtable}

详见 \texttt{08\_theory\_connections.md} 完整对比表。

\begin{center}\rule{0.5\linewidth}{0.5pt}\end{center}

\subsection{理论状态评估}\label{ux7406ux8bbaux72b6ux6001ux8bc4ux4f30}

\begin{longtable}[]{@{}
  >{\raggedright\arraybackslash}p{(\linewidth - 4\tabcolsep) * \real{0.3333}}
  >{\raggedright\arraybackslash}p{(\linewidth - 4\tabcolsep) * \real{0.3333}}
  >{\raggedright\arraybackslash}p{(\linewidth - 4\tabcolsep) * \real{0.3333}}@{}}
\toprule\noalign{}
\begin{minipage}[b]{\linewidth}\raggedright
类别
\end{minipage} & \begin{minipage}[b]{\linewidth}\raggedright
比例
\end{minipage} & \begin{minipage}[b]{\linewidth}\raggedright
说明
\end{minipage} \\
\midrule\noalign{}
\endhead
\bottomrule\noalign{}
\endlastfoot
\textbf{严格理论} & \textasciitilde40\% & Theorem SE-1
的有限状态版本、贝叶斯鞅收敛、Robbins-Monro 标准收敛、Regret 上界 \\
\textbf{形式化猜想} & \textasciitilde20\% &
连续状态空间扩展、双时间尺度收敛、KL 收缩率 \\
\textbf{合理假设} & \textasciitilde40\% & 精确 Lyapunov
函数形式、极限环条件、Gödel 类比、最优退火策略 \\
\end{longtable}

\begin{center}\rule{0.5\linewidth}{0.5pt}\end{center}

\subsection{开放问题（优先级排序）}\label{ux5f00ux653eux95eeux9898ux4f18ux5148ux7ea7ux6392ux5e8f}

\begin{longtable}[]{@{}lll@{}}
\toprule\noalign{}
优先级 & 问题 & 文件 \\
\midrule\noalign{}
\endhead
\bottomrule\noalign{}
\endlastfoot
\textbf{P0} & 明确的 Lyapunov 函数 \(\Phi\) 的定义与下降证明 &
\texttt{06}, \texttt{09} \\
\textbf{P1} & 四种收敛路径的精确定界 & \texttt{06} \\
\textbf{P1} & 耦合 \((S_t, \theta_t)\) 系统的紧收敛率 & \texttt{05},
\texttt{06} \\
\textbf{P2} & 最优 gatekeeper 更新频率 & \texttt{02}, \texttt{03} \\
\textbf{P2} & \(\varepsilon\)-收敛所需的最小记忆库大小 & \texttt{07} \\
\textbf{P3} & 非渐近（有限时间）保证 & \texttt{03}, \texttt{06} \\
\textbf{P4} & 分布偏移下的鲁棒性刻画 & \texttt{05} \\
\textbf{P5} & Gödel 类比的严格化 & \texttt{07} \\
\end{longtable}

\begin{center}\rule{0.5\linewidth}{0.5pt}\end{center}

\subsection{文件结构}\label{ux6587ux4ef6ux7ed3ux6784}

\begin{verbatim}
theory/self_evolution/
├── README.md                              ← 本文件
├── 01_symbol_system.md                    # 符号系统与问题设定
├── 02_dynamical_system.md                 # 离散动力系统形式化
├── 03_online_learning_regret.md           # 在线学习 Regret 分析
├── 04_bayesian_update.md                  # 贝叶斯更新解释
├── 05_stochastic_approximation.md         # 随机逼近分析
├── 06_fixed_point_convergence.md          # 不动点与收敛定理（中心）
├── 07_completeness.md                     # 完备性分析
├── 08_theory_connections.md               # 与已知理论的连接
└── 09_verification_report.md              # 验证报告
\end{verbatim}

\begin{center}\rule{0.5\linewidth}{0.5pt}\end{center}

\subsection{跨领域适用}\label{ux8de8ux9886ux57dfux9002ux7528}

\begin{longtable}[]{@{}
  >{\raggedright\arraybackslash}p{(\linewidth - 6\tabcolsep) * \real{0.1053}}
  >{\raggedright\arraybackslash}p{(\linewidth - 6\tabcolsep) * \real{0.1930}}
  >{\raggedright\arraybackslash}p{(\linewidth - 6\tabcolsep) * \real{0.3333}}
  >{\raggedright\arraybackslash}p{(\linewidth - 6\tabcolsep) * \real{0.3684}}@{}}
\toprule\noalign{}
\begin{minipage}[b]{\linewidth}\raggedright
领域
\end{minipage} & \begin{minipage}[b]{\linewidth}\raggedright
\(S_t\) 角色
\end{minipage} & \begin{minipage}[b]{\linewidth}\raggedright
\(f_{\theta_t}\) 角色
\end{minipage} & \begin{minipage}[b]{\linewidth}\raggedright
\(\mathcal{M}_t\) 角色
\end{minipage} \\
\midrule\noalign{}
\endhead
\bottomrule\noalign{}
\endlastfoot
\textbf{MLIP} & 判断 DFT 标签正确性 & NEP 势能面模型 & 累积的 DFT
计算轨迹 \\
\textbf{医学图像} & 判断诊断标签质量 & 辅助诊断模型 &
累积的专家复核记录 \\
\textbf{LLM} & 判断标注/生成质量 & 下游任务模型 & 累积的人工审核日志 \\
\textbf{自动驾驶} & 判断检测框/标注质量 & 感知模型 &
累积的路测验证数据 \\
\end{longtable}

\begin{center}\rule{0.5\linewidth}{0.5pt}\end{center}

\subsection{参考文献（新增）}\label{ux53c2ux8003ux6587ux732eux65b0ux589e}

\begin{enumerate}
\def\labelenumi{\arabic{enumi}.}
\tightlist
\item
  Robbins, H., \& Monro, S. (1951). A stochastic approximation method.
  \emph{Annals of Mathematical Statistics}, 22(3), 400-407.
\item
  Borkar, V. S. (2008). \emph{Stochastic Approximation: A Dynamical
  Systems Viewpoint}. Cambridge University Press.
\item
  Kushner, H. J., \& Yin, G. G. (2003). \emph{Stochastic Approximation
  and Recursive Algorithms and Applications} (2nd ed.). Springer.
\item
  Doob, J. L. (1953). \emph{Stochastic Processes}. Wiley.
\item
  Bernardo, J. M., \& Smith, A. F. M. (1994). \emph{Bayesian Theory}.
  Wiley.
\item
  van der Vaart, A. W. (1998). \emph{Asymptotic Statistics}. Cambridge
  University Press.
\item
  Silver, D., et al.~(2017). Mastering the game of Go without human
  knowledge. \emph{Nature}, 550, 354-359.
\item
  Shahriari, B., et al.~(2016). Taking the human out of the loop: A
  review of Bayesian optimization. \emph{Proceedings of the IEEE},
  104(1), 148-175.
\item
  Settles, B. (2009). Active learning literature survey.
  \emph{University of Wisconsin-Madison Technical Report}.
\item
  Solomonoff, R. J. (1964). A formal theory of inductive inference.
  \emph{Information and Control}, 7(1), 1-22, 224-254.
\item
  Li, M., \& Vitányi, P. (2008). \emph{An Introduction to Kolmogorov
  Complexity and Its Applications} (3rd ed.). Springer.
\item
  Zinkevich, M. (2003). Online convex programming and generalized
  infinitesimal gradient ascent. \emph{ICML}.
\item
  Cesa-Bianchi, N., \& Lugosi, G. (2006). \emph{Prediction, Learning,
  and Games}. Cambridge University Press.
\item
  SCX Theorem 1-3. \texttt{../theorems/}.
\item
  SCX Unified Theorem Document. \texttt{../THEOREMS\_UNIFIED.md}.
\end{enumerate}

\begin{center}\rule{0.5\linewidth}{0.5pt}\end{center}

\emph{End of README.md --- SCX Self-Evolution Theory}

\end{document}
